\section{Task, models, and set-size manipulation}

We study a recurrent, pixel-based agent performing a spatially-cued orientation-change detection task and ask how the cost of an \emph{invalid} cue scales with the number of competing locations. The task is value-directed covert attention \citep{luo2019attention,carrasco2011visual}: on each trial a coloured cue marks one location and, by its colour, signals an associated reward magnitude (the cue colour-values used were $[1,3,5]$). The cue's spatial reliability is set by a validity proportion. Each episode unfolds over seven logical frames; on half of trials an oriented stimulus changes its orientation, and when a change occurs it appears at frame~5 (f5), so the change is detectable from f5 onward, with f6 being one frame later. The agent emits a \texttt{wait}/\texttt{change} decision on every frame. A \textbf{valid} trial is one in which the change occurs at the cued location; an \textbf{invalid} trial is one in which the change occurs at an uncued location, with the cue therefore pointing away from the event \citep{posner1980attention,posner1994visual}.

The agent is a convolutional JEPA architecture trained end-to-end from raw pixels. A squeeze-and-excitation residual convolutional front-end (no variational encoder) extracts per-location features; a V-JEPA prototype-softmax self-distillation objective provides an auxiliary representational signal; a cross-attention memory module (the cue-gated, top-down ``crossattn1'' reader) feeds a spatial xLSTM recurrent memory (memory width $d_{\mathrm{mem}}=128$); and a distributional actor-critic head produces the decision. The cross-attention reader is the priority channel of interest: its queries read out over the spatial memory keys, so the distribution of attention over those keys reports which location the agent is monitoring on each frame.

We manipulate stimulus load by changing only the set size, holding the architecture, auxiliary objective, training regime, and decision head fixed. In the \textbf{set size 4} condition (SS4, task vda4) the four stimuli occupy a $2\times2$ grid ($25\times25$-pixel patches in a $50\times50$ image). In the \textbf{set size 9} condition (SS9, task vda9) nine stimuli occupy a $3\times3$ grid ($75\times75$ image), with the cue placed at the top-left location (P0) or the bottom-right location (P8). Because the same model and objective are used at both set sizes, any difference in how the invalid cue is handled is attributable to load rather than to a change in mechanism. As a mechanistic control we additionally examine an element-wise-affine variant of the same conv-JEPA at SS9, in which the memory reader is a bottom-up self-attention rather than the cue-gated cross-attention; this isolates the contribution of the top-down channel.
